% =====================================================================
%  06_prototipos.tex — Un apartado por prototipo, con la misma
%  estructura: introducción, objetivo, historias de usuario,
%  arquitectura, componentes, interfaz, desarrollo y pruebas.
%  Al agregar el Prototipo 2, copiar la estructura de la sección 1.
% =====================================================================
\chapter{Prototipos}
\label{cap:prototipos}

\section{Prototipo 1: Sistema base en sitio}
\label{sec:prototipo1}

\subsection{Introducción}

El Prototipo~1 constituye la base funcional de \saru{}: implementa la
cadena completa desde la cámara hasta el registro y el aviso, sobre la que
se construirán los prototipos siguientes. Resulta de la fusión, el 3 de
noviembre de 2026, de dos versiones previas del sistema (la del repositorio
y la que corría en la Jetson), seguida de correcciones con video real y de
su puesta en operación en el patio~\cite{rep-fusion,rep-envivo}. Establece
la arquitectura en capas, el esquema de la base de datos, los tres modelos
en cascada y la lógica de cruce que se reutilizarán en las siguientes
iteraciones.

\subsection{Objetivo}

Implementar un prototipo funcional que, a partir del video de la cámara de
la puerta del Patio CCA, detecte cada unidad, lea su número económico,
determine si entra o sale, registre el movimiento con evidencia y lo
notifique, y que ofrezca un tablero web para consultar los movimientos y
las grabaciones.

\subsection{Especificaciones (historias de usuario) del Prototipo 1}

Las historias se escriben desde el punto de vista de cuatro tipos de
usuario: \textbf{operador del patio}, \textbf{validador},
\textbf{coordinación} y \textbf{personal técnico}.

\begin{aviso}[title={Por validar}]
Las historias se derivaron del sistema existente y de los reportes del
equipo. Deben confirmarse con los usuarios reales del patio.
\end{aviso}

\historiaUsuario{HU-01}{Consultar unidades dentro}
  {Como operador del patio, quiero ver qué unidades están dentro y desde
   cuándo, para saber de qué unidades dispongo.}
  {Alta}
  {\begin{itemize}[nosep,leftmargin=*]
     \item El tablero muestra las unidades con estancia abierta y los
           minutos que llevan dentro.
     \item La lista se actualiza sin recargar la página.
   \end{itemize}}

\historiaUsuario{HU-02}{Ver video en vivo}
  {Como operador del patio, quiero ver el video de la puerta en vivo, para
   comprobar lo que la cámara está viendo.}
  {Alta}
  {\begin{itemize}[nosep,leftmargin=*]
     \item El video se muestra a 10~cuadros por segundo y 960~px de ancho.
     \item Si se pierde la cámara, el tablero indica ``Sin señal''.
   \end{itemize}}

\historiaUsuario{HU-03}{Consultar cruces}
  {Como operador del patio, quiero consultar los cruces registrados y
   buscarlos por número, empresa o fecha, para revisar un movimiento
   concreto.}
  {Alta}
  {\begin{itemize}[nosep,leftmargin=*]
     \item La lista muestra número, empresa, dirección, zona, hora y foto.
     \item Permite buscar por número o empresa y filtrar por fechas.
     \item Muestra 10 cruces por página, con navegación entre páginas.
   \end{itemize}}

\historiaUsuario{HU-04}{Recibir aviso de cruce}
  {Como operador del patio, quiero recibir un aviso con fotos cada vez que
   una unidad entra o sale, para enterarme sin estar frente al tablero.}
  {Alta}
  {\begin{itemize}[nosep,leftmargin=*]
     \item El aviso indica número, dirección, tipo y empresa.
     \item Incluye hasta 3 fotos de momentos distintos, separadas al menos
           1~segundo.
     \item Se envía por Telegram y, si está configurado, por WhatsApp.
   \end{itemize}}

\historiaUsuario{HU-05}{Validar una lectura dudosa}
  {Como validador, quiero confirmar o corregir el número de un cruce
   dudoso desde el aviso, para que el registro quede correcto.}
  {Media}
  {\begin{itemize}[nosep,leftmargin=*]
     \item El aviso incluye botones para confirmar o corregir.
     \item Solo los usuarios autorizados pueden validar.
     \item Antes del aviso, Gemini puede revisar la foto si está activado.
   \end{itemize}}

\historiaUsuario{HU-06}{Registrar cada cruce una sola vez}
  {Como coordinación, quiero que cada paso de una unidad genere un solo
   registro con el número correcto, para que los conteos sean confiables.}
  {Alta}
  {\begin{itemize}[nosep,leftmargin=*]
     \item Un cruce se registra cuando la unidad termina de pasar, con el
           número más leído y las mejores fotos.
     \item Una unidad detenida sobre la línea no genera registros.
     \item Una maniobra (regreso antes de 20~s) no genera registro.
     \item El mismo número no se registra dos veces en menos de 60~s.
     \item Con varias unidades a la vez, cada una lleva su propio cruce.
   \end{itemize}}

\historiaUsuario{HU-07}{Registrar números fuera del padrón}
  {Como coordinación, quiero que un número leído de forma consistente se
   registre aunque no esté en el padrón, para no perder movimientos ni
   inventar unidades.}
  {Media}
  {\begin{itemize}[nosep,leftmargin=*]
     \item Con 4 lecturas iguales, el número se registra tal cual, marcado
           como fuera del padrón.
     \item No se corrige a un número parecido del padrón.
     \item El tablero permite consultar los números fuera del padrón.
   \end{itemize}}

\historiaUsuario{HU-08}{Consultar estancias y trayectoria}
  {Como coordinación, quiero consultar las estancias y el recorrido de una
   unidad, para conocer cuánto tiempo permanece en el patio.}
  {Media}
  {\begin{itemize}[nosep,leftmargin=*]
     \item Cada estancia muestra hora de entrada, hora de salida y
           duración en minutos.
     \item El día operativo empieza a las 3:00~a.m.
     \item Se muestra el promedio, mínimo y máximo de estancia por
           empresa.
   \end{itemize}}

\historiaUsuario{HU-09}{Consultar grabaciones}
  {Como coordinación, quiero reproducir el video de una hora concreta,
   para revisar un cruce cuestionado.}
  {Alta}
  {\begin{itemize}[nosep,leftmargin=*]
     \item Hay un archivo por hora de las últimas 72~horas.
     \item El tablero lista las grabaciones por cámara y permite
           reproducirlas.
     \item La grabación continúa aunque falle la detección.
   \end{itemize}}

\historiaUsuario{HU-10}{Revisar el estado del sistema}
  {Como personal técnico, quiero conocer el estado del sistema, para
   detectar fallas sin entrar a la Jetson.}
  {Alta}
  {\begin{itemize}[nosep,leftmargin=*]
     \item Una ruta de salud informa conexión de la cámara, cuadros por
           segundo y contadores.
     \item Si la cámara se desconecta, el sistema intenta reconectarse cada
           2~segundos.
   \end{itemize}}

\historiaUsuario{HU-11}{Probar con video grabado}
  {Como personal técnico, quiero ejecutar el sistema con un video en lugar
   de la cámara, para probar cambios sin afectar la operación.}
  {Media}
  {\begin{itemize}[nosep,leftmargin=*]
     \item Con un archivo de video como fuente, el sistema funciona igual
           que con la cámara.
     \item El tablero permite avanzar, retroceder y saltar entre videos.
     \item Los avisos pueden desactivarse durante las pruebas.
   \end{itemize}}

\historiaUsuario{HU-12}{Actualizar el padrón}
  {Como personal técnico, quiero dar de alta o de baja unidades editando un
   archivo, para mantener el padrón al día sin tocar la base de datos.}
  {Media}
  {\begin{itemize}[nosep,leftmargin=*]
     \item El padrón se lee de un archivo CSV en cada arranque.
     \item Si una fila es inválida, el sistema no arranca e indica la
           línea con el error.
   \end{itemize}}

\historiaUsuario{HU-13}{Reunir recortes para entrenamiento}
  {Como personal técnico, quiero que el sistema guarde recortes de las
   unidades, para ampliar los datos con que se entrenan los modelos.}
  {Media}
  {\begin{itemize}[nosep,leftmargin=*]
     \item Se guarda un recorte por unidad cada cierto número de segundos.
     \item El tablero muestra estadísticas de los recortes guardados.
   \end{itemize}}

\subsection{Arquitectura del Prototipo 1}

El Prototipo~1 aplica la arquitectura en capas del
capítulo~\ref{cap:bosquejo}. Un solo proceso de Python
(\texttt{cisa-vision}) arranca la cámara, lanza dos hilos de visión y el
servidor web. La cámara publica el cuadro más reciente, del que leen la
detección, el video en vivo y las capturas. El primer hilo ejecuta M1 y el
seguimiento; el segundo, M2 y M3 sobre los recortes. Cuando un número
confirmado cruza la línea, el sistema espera a que la unidad termine de
pasar, valida, registra y avisa. La grabación lee la cámara por separado.

\begin{figure}[htbp]
  \centering
  \resizebox{\linewidth}{!}{%
  \begin{tikzpicture}
    \node[caja] (cam)   at (0,0)    {Cámara\\RTSP};
    \node[caja] (m1)    at (3.8,0)  {M1\\YOLO + ByteTrack};
    \node[caja] (m23)   at (7.6,0)  {M2 + M3\\zona y dígitos};
    \node[caja] (voto)  at (11.2,0) {Votación\\lecturas iguales};
    \node[caja] (linea) at (11.2,-2.1) {Línea virtual\\entrada / salida};
    \node[caja] (valid) at (7.6,-2.1)  {Validación\\Gemini + Telegram};
    \node[caja] (reg)   at (3.8,-2.1)  {Registro\\SQLite};
    \node[caja] (web)   at (0,-2.1)    {Tablero web\\MJPEG + consultas};
    \node[cajaGris] (grab)   at (0,-4.2)   {Grabación\\lee la cámara};
    \node[cajaGris] (avisos) at (3.8,-4.2) {Avisos\\Telegram / WhatsApp};
    \draw[flecha] (cam) -- (m1);
    \draw[flecha] (m1) -- node[etiqueta,above]{recortes} (m23);
    \draw[flecha] (m23) -- node[etiqueta,above]{número} (voto);
    \draw[flecha] (voto) -- (linea);
    \draw[flecha] (linea) -- node[etiqueta,above]{cruce} (valid);
    \draw[flecha] (valid) -- (reg);
    \draw[flecha] (reg) -- (web);
    \draw[flecha] (reg) -- (avisos);
    \draw[flechaGris] (cam.west) -- ++(-0.6,0) |- (grab.west);
    \begin{scope}[on background layer]
      \node[draw=azulCISA, dashed, rounded corners=4pt,
            fit=(m1)(m23)(voto)(linea)(valid)(reg), inner sep=0.3cm,
            label={[font=\scriptsize\bfseries,text=azulCISA]above:%
                   {cisa-vision (un proceso, dos hilos)}}] {};
    \end{scope}
  \end{tikzpicture}%
  }
  \caption{Arquitectura del Prototipo 1: flujo de una unidad desde la
           cámara hasta el registro.}
  \label{fig:arq-p1}
\end{figure}

\subsection{Componentes del Prototipo 1}

\begin{longtable}{@{}L{4.4cm}L{2.8cm}L{6.8cm}@{}}
  \caption{Componentes del Prototipo 1.}
  \label{tab:componentes-p1} \\
  \toprule
  \textbf{Componente} & \textbf{Capa} & \textbf{Responsabilidad} \\
  \midrule
  \endfirsthead
  \caption[]{Componentes del Prototipo 1 (continuación).} \\
  \toprule
  \textbf{Componente} & \textbf{Capa} & \textbf{Responsabilidad} \\
  \midrule
  \endhead
  \path{config.py} & Configuración & Único lector de la configuración;
    valida variables, modelos, carpetas y padrón al arrancar. \\
  \path{pipe_stream_v2.py} & Arranque & Crea la cámara, lanza los hilos y
    el servidor web. \\
  \path{logica/camara.py} & Visión & Lectura RTSP, respaldo con FFmpeg,
    reconexión y modo prueba. \\
  \path{logica/estado.py} & Visión & Cuadro más reciente y estado
    compartido del sistema. \\
  \path{ml/modelos.py} & Visión & Carga y calentamiento de los motores
    TensorRT. \\
  \path{ml/hilos.py} & Visión & Hilo de detección (M1) e hilo de lectura
    (M2 y M3); registro único por cruce. \\
  \path{ml/reconocimiento.py} & Visión & Recortes, lectura de dígitos,
    validación del número y guardado de recortes. \\
  \path{ml/linea_virtual.py} & Visión & Lado de la línea, margen,
    maniobras y evaluación del cruce. \\
  \path{metricas/schema.sql} & Datos & Tablas, índices y vistas. \\
  \path{metricas/storage.py} & Datos & Único acceso a la base;
    sincronización del padrón, altas y consultas. \\
  \path{avisos/telegram.py} & Servicios externos & Avisos y validación
    por Telegram. \\
  \path{avisos/gemini.py} & Servicios externos & Revisión automática de
    la foto. \\
  \path{avisos/whatsapp.py}, \path{whatsapp/} & Servicios externos &
    Puente hacia WhatsApp Web (Node.js y Baileys). \\
  \path{web/rutas.py} & Presentación & Rutas del tablero, API, video y
    archivos. \\
  \path{templates/}, \path{static/} & Presentación & Páginas, estilos y
    JavaScript del tablero. \\
  \path{scripts/grabar.sh} & Datos & Grabación continua y rotación de
    archivos. \\
  \bottomrule
\end{longtable}

\subsection{Diseño de la interfaz del Prototipo 1}

La interfaz del Prototipo~1 tiene tres páginas: el tablero principal, el
video en vivo con la lista de cruces y la consulta de grabaciones. Además,
los avisos llegan a Telegram y WhatsApp.

\figuraPendiente{06\_prototipos/p1-tablero.png}
  {Tablero principal: indicadores, gráficas y unidades dentro}
  {fig:p1-tablero}

\figuraPendiente{06\_prototipos/p1-livevideo.png}
  {Video en vivo y lista paginada de cruces}
  {fig:p1-livevideo}

\figuraPendiente{06\_prototipos/p1-grabaciones.png}
  {Consulta y reproducción de grabaciones por hora}
  {fig:p1-grabaciones}

\figuraPendiente{06\_prototipos/p1-aviso.png}
  {Aviso de cruce en Telegram con fotos y botones de validación}
  {fig:p1-aviso}

\subsection{Desarrollo del Prototipo 1}

El Prototipo~1 se consolidó entre el 3 y el 12 de noviembre de 2026. La
tabla~\ref{tab:desarrollo-p1} resume los cambios principales según los
reportes del equipo.

\begin{longtable}{@{}L{2.2cm}L{11.8cm}@{}}
  \caption{Cronología del desarrollo del Prototipo 1.}
  \label{tab:desarrollo-p1} \\
  \toprule
  \textbf{Fecha} & \textbf{Cambio} \\
  \midrule
  \endfirsthead
  \toprule
  \textbf{Fecha} & \textbf{Cambio} \\
  \midrule
  \endhead
  3 nov & Fusión de las dos versiones previas: varias unidades a la vez y
    topes de procesamiento (repositorio) con margen de línea, espera de
    25~s, avisos por WhatsApp y padrón de 396 unidades (Jetson). La versión
    fusionada pasa 16 de 16 pruebas del simulador, y la detección sube de
    13.4 a 22.1 cuadros por segundo en modo de 40~W~\cite{rep-fusion}. \\
  4 nov & Correcciones con video real: tres fallas que impedían registrar
    unidades bien leídas; los números fuera del padrón ahora se registran y
    se muestran en el tablero. \\
  5--9 nov & Sistema en vivo en el patio: tres días sin fallas. Los avisos
    de Telegram no salen porque la red del patio bloquea el servicio; se
    preparó la solicitud para habilitarlo~\cite{rep-envivo}. \\
  9--10 nov & Panel de cruces en la página de video en vivo, con búsqueda,
    filtros y paginación de 10. \\
  10 nov & Incidente: avisos de WhatsApp que llegan sin poder leerse en el
    teléfono destinatario~\cite{rep-whatsapp}. \\
  11 nov & Hasta 3 fotos por aviso; visor de fotos; videos de prueba del
    más reciente al más antiguo. \\
  12 nov & Línea virtual de varios puntos que termina en sus extremos y
    compara la base del recuadro de la unidad. Probada en simulación;
    pendiente con video real. Solo en el paquete
    \texttt{SARU\_V4}~\cite{rep-tecnico}. \\
  \bottomrule
\end{longtable}

\begin{aviso}[title={Versiones del código}]
Existen al menos tres versiones del Prototipo~1: la del repositorio, el
paquete \texttt{SARU\_V4} (con la línea de varios puntos) y la copia de la
Jetson. Este documento describe la del repositorio salvo que se indique lo
contrario. Definir una sola fuente de verdad es parte del trabajo a futuro
(capítulo~\ref{cap:futuro}).
\end{aviso}

\figuraPendiente{06\_prototipos/p1-secuencia-cruce.png}
  {Diagrama de secuencia del registro de un cruce: detección, lectura,
   votación, cruce, espera, validación, registro y aviso}
  {fig:p1-secuencia}

\subsection{Protocolo de pruebas del Prototipo 1}
\label{sec:pruebas-p1}

Cada caso de prueba se identifica con una sigla de módulo y un número
consecutivo; por ejemplo, \texttt{LVC-01} es el primer caso de las pruebas
de línea virtual y cruce. La columna \textbf{Aprobado} indica
\textbf{Sí}, \textbf{No} o \textbf{Pendiente} (sin evidencia todavía). La
evidencia de cada caso se reúne en el apéndice~\ref{cap:evidencias-p1}.

\begin{longtable}{@{}L{1.4cm}L{4.0cm}L{8.6cm}@{}}
  \caption{Nomenclatura de las pruebas del Prototipo 1.}
  \label{tab:nomenclatura-p1} \\
  \toprule
  \textbf{Sigla} & \textbf{Módulo} & \textbf{Qué verifica} \\
  \midrule
  \endfirsthead
  \toprule
  \textbf{Sigla} & \textbf{Módulo} & \textbf{Qué verifica} \\
  \midrule
  \endhead
  ADQ & Adquisición de video & Conexión, reconexión y modo prueba. \\
  DLN & Detección y lectura & Detección de unidades, lectura y
    confirmación del número. \\
  LVC & Línea virtual y cruce & Dirección del cruce, margen, maniobras,
    varias unidades y registro único. \\
  REG & Registro & Eventos, estancias, padrón y números fuera del padrón. \\
  AVS & Avisos y validación & Entrega de avisos y validación. \\
  WEB & Tablero web & Páginas, consultas y video en vivo. \\
  GRB & Grabación & Archivos por hora, rotación y consulta. \\
  REN & Rendimiento & Cuadros por segundo, consumo y estabilidad. \\
  \bottomrule
\end{longtable}

\subsubsection{Módulo de línea virtual y cruce (LVC)}

\textbf{Objetivo:} verificar que el sistema registra una sola entrada o
salida por cada paso real, y ninguna cuando la unidad se detiene, maniobra
o rodea la línea.

\textbf{Componentes que se prueban:} \path{ml/linea_virtual.py} y el
registro único de \path{ml/hilos.py}.

\textbf{Entorno:} simulador del equipo que reproduce cruces con reloj y
modelos falsos~\cite{rep-fusion,rep-tecnico}. No incluye los modelos YOLO
ni video real. El simulador no está en el repositorio.

{\small
\begin{longtable}{@{}L{1.3cm}L{4.3cm}L{3.8cm}C{1.4cm}L{3.2cm}@{}}
  \caption{Cobertura de pruebas del módulo LVC.}
  \label{tab:lvc} \\
  \toprule
  \textbf{ID} & \textbf{Caso} & \textbf{Resultado esperado} &
    \textbf{Aprobado} & \textbf{Observaciones} \\
  \midrule
  \endfirsthead
  \toprule
  \textbf{ID} & \textbf{Caso} & \textbf{Resultado esperado} &
    \textbf{Aprobado} & \textbf{Observaciones} \\
  \midrule
  \endhead
  LVC-01 & Una unidad entra. & Una entrada. & Sí & Simulador. \\
  LVC-02 & Una unidad sale. & Una salida. & Sí & Simulador. \\
  LVC-03 & Paso rápido (2~s) y lento (40~s). & Un registro en cada caso. &
    Sí & Simulador. \\
  LVC-04 & La unidad regresa a los 5~s. & Maniobra, sin registro. & Sí &
    Simulador. \\
  LVC-05 & La unidad regresa a los 15~s. & Maniobra, sin registro. & Sí &
    Fallaba en las versiones previas. \\
  LVC-06 & Unidad detenida sobre la línea con temblor. & Sin registro. & Sí
    & 20 de 20 corridas. \\
  LVC-07 & Unidad detenida con vaivén lento. & Sin registro. & Sí & 20 de
    20 corridas. \\
  LVC-08 & Dos unidades lado a lado. & Un registro por unidad. & Sí &
    Fallaba en la versión de la Jetson. \\
  LVC-09 & Una unidad estacionada y otra entra. & Solo la entrada de la
    segunda. & Sí & Simulador. \\
  LVC-10 & Tres unidades a la vista, solo una cruza. & Un solo registro. &
    Sí & Simulador. \\
  LVC-11 & Una entra y otra sale a la vez. & Una entrada y una salida. & Sí
    & Simulador. \\
  LVC-12 & La unidad rodea la línea por un extremo. & Sin registro. & Sí &
    Solo línea de varios puntos (\texttt{SARU\_V4}). \\
  LVC-13 & Entrada y salida con la línea de varios puntos y video real. &
    Un registro por paso. & Pendiente & Validar en la Jetson. \\
  \bottomrule
\end{longtable}
}

\subsubsection{Módulos de registro, avisos y rendimiento (REG, AVS, REN)}

\textbf{Objetivo:} verificar que los cruces quedan registrados, que los
avisos llegan y que el sistema opera de forma estable en la Jetson.

{\small
\begin{longtable}{@{}L{1.3cm}L{4.3cm}L{3.8cm}C{1.4cm}L{3.2cm}@{}}
  \caption{Cobertura de pruebas de los módulos REG, AVS y REN.}
  \label{tab:reg-avs-ren} \\
  \toprule
  \textbf{ID} & \textbf{Caso} & \textbf{Resultado esperado} &
    \textbf{Aprobado} & \textbf{Observaciones} \\
  \midrule
  \endfirsthead
  \toprule
  \textbf{ID} & \textbf{Caso} & \textbf{Resultado esperado} &
    \textbf{Aprobado} & \textbf{Observaciones} \\
  \midrule
  \endhead
  REG-01 & Unidad fuera del padrón leída de forma consistente. & Se
    registra con su número y aparece en el tablero. & Sí & Video real,
    4~nov. \\
  REG-02 & Memoria tras 300 unidades. & Ningún estado sin uso. & Sí &
    Simulador: 0 registros sin uso. \\
  REG-03 & Entrada y salida de la misma unidad. & Estancia abierta y luego
    cerrada con su duración. & Pendiente & --- \\
  AVS-01 & Aviso por Telegram desde la red del patio. & El aviso llega al
    grupo. & No & La red del patio bloquea Telegram~\cite{rep-envivo}. \\
  AVS-02 & Aviso por WhatsApp al destinatario principal. & El mensaje
    llega y se puede leer. & No & Mensajes ilegibles el
    10~nov~\cite{rep-whatsapp}. \\
  AVS-03 & Aviso por WhatsApp a otro número. & El mensaje llega y se puede
    leer. & Sí & Prueba de aislamiento del incidente. \\
  REN-01 & Detección en modo de 40~W. & Más de 20~cuadros por segundo. &
    Sí & 22.1~cuadros por segundo~\cite{rep-fusion}. \\
  REN-02 & Operación continua en el patio. & Sin fallas durante varios
    días. & Sí & 3 días; carga y temperatura normales~\cite{rep-envivo}. \\
  REN-03 & Consumo en modo de 25~W. & Bajo consumo. & Sí & 11~W sin unidad;
    14 a 19~W con unidad~\cite{rep-tecnico}. \\
  \bottomrule
\end{longtable}
}

\subsubsection{Módulos pendientes de prueba (ADQ, DLN, WEB, GRB)}

{\small
\begin{longtable}{@{}L{1.3cm}L{4.3cm}L{3.8cm}C{1.4cm}L{3.2cm}@{}}
  \caption{Cobertura de pruebas de los módulos ADQ, DLN, WEB y GRB.}
  \label{tab:adq-dln-web-grb} \\
  \toprule
  \textbf{ID} & \textbf{Caso} & \textbf{Resultado esperado} &
    \textbf{Aprobado} & \textbf{Observaciones} \\
  \midrule
  \endfirsthead
  \toprule
  \textbf{ID} & \textbf{Caso} & \textbf{Resultado esperado} &
    \textbf{Aprobado} & \textbf{Observaciones} \\
  \midrule
  \endhead
  ADQ-01 & Se desconecta la cámara. & El sistema reintenta cada 2~s y se
    recupera solo. & Pendiente & Implementado en el código. \\
  ADQ-02 & Fuente de video de prueba. & Funciona igual que la cámara. &
    Pendiente & --- \\
  DLN-01 & Lectura con vista lateral. & Número confirmado correcto. &
    Pendiente & Falta medir la exactitud. \\
  DLN-02 & Lectura con vista delantera o trasera. & Número confirmado
    correcto. & No & Según pruebas del equipo, M2 casi solo encuentra la
    zona con vista lateral. \\
  WEB-01 & Búsqueda y paginación de cruces. & Resultados filtrados, 10 por
    página. & Pendiente & --- \\
  WEB-02 & Video en vivo. & 10~cuadros por segundo, 960~px de ancho. &
    Pendiente & --- \\
  GRB-01 & Grabación de 72~h. & Un archivo por hora; se borran los de más
    de 72~h. & Pendiente & --- \\
  GRB-02 & Falla la detección. & La grabación continúa. & Pendiente &
    --- \\
  \bottomrule
\end{longtable}
}

\subsection{Conclusiones del Prototipo 1}

El Prototipo~1 cumple la cadena completa de detección, lectura, cruce,
registro, avisos, tablero y grabación, y opera en vivo en el patio. La
lógica de cruce es el módulo con más evidencia: pasa todos los casos del
simulador, incluidos los de varias unidades a la vez que fallaban en las
versiones previas. Los puntos débiles son la \textbf{entrega de avisos},
que hoy falla por la red del patio y por el incidente de WhatsApp, y la
\textbf{lectura en vistas no laterales}. Además, falta evidencia formal de
los módulos de adquisición, tablero y grabación, y medir la exactitud de la
lectura con una muestra revisada contra el video.
